\documentclass[12pt]{article}
\usepackage[T1]{fontenc}
\usepackage[utf8]{inputenc}
\usepackage{lmodern}
\usepackage{textcomp}
\usepackage{enumitem}
\usepackage{geometry}
\geometry{margin=1in}
\begin{document}
\begin{center}
Answer Key - Political Science - Undergraduate Level Assessment 2 \
\small (Answers are ordered exactly as in the question paper.)
\end{center}

\section*{Multiple Choice}
\begin{enumerate}[label=\arabic*.]
\item \textbf{Answer:} B
\\ \textit{Options:} A) Fewer than 7; B) Between 8 and 15; C) Between 16 and 25; D) More than 25
\\ \textit{Note:} The correct answer is between 8 and 15 because, as of now, there are officially recognized nuclear-armed states and several others that are suspected of possessing nuclear capabilities, reflecting the complex dynamics of nuclear proliferation in the international system.
\item \textbf{Answer:} C
\\ \textit{Options:} A) resigns to take a position as a consultant to lobbying groups; B) cedes control of his or her department to lifelong bureaucrats within the department; C) places his or her department's priorities above the president's; D) accepts bribes or expensive favors from businesses regulated by his or her department
\\ \textit{Note:} The phrase "gone native" refers to a cabinet member prioritizing their department's interests over the president's agenda, which can lead to conflicts and undermine the unified direction of the administration.
\item \textbf{Answer:} C
\\ \textit{Options:} A) Deploying troops; B) Drafting appropriations bills; C) Negotiating treaties; D) Forcing Congress into session
\\ \textit{Note:} Negotiating treaties is a power shared with the Senate because the U.S. Constitution requires that treaties be ratified by a two-thirds majority in the Senate after the President negotiates them.
\item \textbf{Answer:} C
\\ \textit{Options:} A) The right to privacy is determined entirely by the states on a case-by-case basis.; B) The right to privacy is explicitly granted in the Preamble to the Constitution.; C) The Supreme Court has ruled that the right to privacy is implied by the Bill of Rights.; D) Common law requires the government to respect citizens' right to privacy.
\\ \textit{Note:} The correct answer is accurate because the Supreme Court has interpreted various amendments in the Bill of Rights to collectively imply a right to privacy, establishing that this right is not explicitly mentioned but is essential to the protection of individual liberties.
\item \textbf{Answer:} C
\\ \textit{Options:} A) Societal groups have a right to survive.; B) Societal groups have their own reality.; C) Societal groups are constituted by social interaction.; D) Societal groups are multiple-identity units.
\\ \textit{Note:} The statement "Societal groups are constituted by social interaction" is the odd one out because it emphasizes the foundational role of interpersonal relationships in forming societal structures, while the other statements likely pertain to distinct aspects of political science that do not focus on social interaction as their primary theme.
\end{enumerate}

\section*{True / False}
\begin{enumerate}[label=\arabic*.]
\setcounter{enumi}{5}
\item \textbf{Answer:} True
\\ \textit{Options:} A) True; B) False
\\ \textit{Note:} American exceptionalism, characterized by a belief in the unique qualities of the United States and its emphasis on individualism and limited government, often fosters a skepticism toward centralized world governance, as it is perceived as a threat to national sovereignty and personal liberties.
\item \textbf{Answer:} True
\\ \textit{Options:} A) True; B) False
\\ \textit{Note:} Leaders may opt to deceive about their nuclear capabilities as a strategic maneuver to maintain a deterrent effect, secure national interests, and mitigate the risk of preemptive strikes from other nations.
\item \textbf{Answer:} False
\\ \textit{Options:} A) True; B) False
\\ \textit{Note:} George H.W. Bush opted not to remove Saddam Hussein from power after the Gulf War primarily to avoid a protracted conflict and the potential destabilization of Iraq, prioritizing regional stability over the expansion of U.S. influence.
\item \textbf{Answer:} True
\\ \textit{Options:} A) True; B) False
\\ \textit{Note:} Greedy states often pursue self-interest driven by human nature, the pursuit of wealth through economic gain, and the promotion of ideologies that justify their actions, reflecting a combination of these motivations in their political behavior.
\item \textbf{Answer:} True
\\ \textit{Options:} A) True; B) False
\\ \textit{Note:} The 2008 financial crisis exposed the vulnerabilities and inequalities inherent in the US model of political economy and capitalism, leading to widespread criticism and a loss of faith in these systems globally.
\end{enumerate}

\section*{Long Form Response}
\begin{enumerate}[label=\arabic*.]
\setcounter{enumi}{10}
\item \textbf{Answer:} The energy security nexus significantly influences international relations by shaping geopolitical rivalries and fostering cooperation among core powers. Control over energy reserves, particularly in resource-rich regions such as the South, has become a central point of contention, leading to heightened tensions among nations vying for dominance. This struggle for energy resources not only affects bilateral relations but also complicates multilateral efforts to address global challenges, such as climate change and security threats. Consequently, the pursuit of energy security can undermine existing alliances and create new rivalries, ultimately impacting international stability and cooperation. As core powers navigate these dynamics, the sustainability of their partnerships often hinges on their ability to balance competition with collaborative engagement in energy governance.
\\ \textit{Note:} correct because it effectively highlights how the competition for energy resources can lead to geopolitical rivalries while also acknowledging the potential for cooperation among core powers, thereby illustrating the dual impact of the energy security nexus on international relations.
\item \textbf{Answer:} Achieving effective international arms control presents numerous challenges, including disparities in national interests, varying levels of military capability, and the complexities of the global defense trade. However, it is possible to overcome these obstacles through strategic cooperation and compromise among governments and various stakeholders, including international organizations and civil society. Building trust through diplomatic engagement and transparency can facilitate collaborative frameworks that address security concerns while promoting disarmament. Additionally, establishing verifiable agreements and monitoring mechanisms can enhance compliance and accountability, thereby strengthening the arms control regime. Ultimately, sustained commitment to dialogue and partnership is essential for fostering a more secure and cooperative international environment.
\\ \textit{Note:} The answer correctly identifies the multifaceted challenges of international arms control while emphasizing the necessity of strategic cooperation and trust-building among governments and stakeholders to address these issues effectively.
\end{enumerate}

\vfill
\noindent\textit{End of Answer Key}
\end{document}